\chapter{抽象方法体系}
本章仅以抽象对象说明小节接口检查所需的结构形态，不表达任何真实研究结论。

\section{前置父节}
本节建立一个位于目标序列之前的抽象单元，用于形成可检查的文档顺序边界。

\subsection{前置单元}
本小节描述占位对象的准备过程，并明确其内容只服务于后续接口测试的结构安排。

该过程在当前范围内完成收束，同时为后续单元提供输入并保留清晰的交接位置。

\section{连续缺口父节}
这里给出另一组彼此独立的抽象描述，父节导语不预先替代各小节的承接与交棒责任。

\subsection{缺口单元甲}
离散容器保存局部符号及其顺序关系，当前描述聚焦容器内部的静态组织方式。

局部符号按照预设次序进入容器，描述在此结束且没有说明下一单元需要接收什么。

\subsection{缺口单元乙}
另一种载体记录临时标签及其组合方式，当前描述聚焦载体内部的排列规则。

临时标签依照独立规则完成排列，描述在此结束且没有指出相邻单元之间的联系。

\subsection{缺口单元丙}
第三类框架整理抽象节点及其局部属性，当前描述聚焦节点集合的静态边界。

抽象节点依据内部约束完成整理，描述在此结束且没有交代后续单元所需的输入。

\section{跨父节后继}
本节引入新的父级范围，并用一段合格导语承担跨父节交接时的只读证据角色。

\subsection{后继单元}
独立索引保存另一批抽象条目及其检索顺序，当前描述聚焦索引内部的组织规则。

本小节在该父节中承担后继示例角色，并以定位句结束这一组跨父节结构。

\section{边界形态}
本节集中放置单段、短段与列表结尾三种边界形态，用于验证合格段筛选规则。

\subsection{单段小节}
本小节只包含一个合格正文段，该段同时承担入口定位与当前单元收束两种结构角色。

\subsection{短段小节}
短段不足二十字。

\subsection{列表结尾小节}
本小节先给出一段长度足够的正文，但该段随后直接连接列表环境并因此属于环境结尾边界。
\begin{itemize}
  \item 抽象条目甲
  \item 抽象条目乙
\end{itemize}
